\chapter{ESPECIFICACIONES TECNICAS DE SUMINISTRO DE MATERIALES}

\section{Alcance}

Las presentes especificaciones describen los materiales principales para la instalación eléctrica interior de la vivienda. La selección final debe cumplir el CNE-U, la EM.010 y los criterios técnicos del proyecto académico.

\section{Conductores eléctricos}

\begin{itemize}
  \item Material: cobre.
  \item Aislamiento: adecuado para instalación interior en tuberia o canalizacion.
  \item Secciones: segun cálculo de corriente, caída de tensión y protección.
  \item Identificacion: fase, neutro y puesta a tierra diferenciados por color cuando sea posible.
\end{itemize}

\begin{landscape}
\begin{table}[H]
\centering
\small
\caption{Conductores propuestos}
\label{tab:conductores-propuestos}
\begin{tabularx}{\textwidth}{L{3.4cm}L{3.3cm}Y}
\toprule
\textbf{Uso} & \textbf{Sección propuesta} & \textbf{Observación} \\
\midrule
Alumbrado & 1.5 mm2 Cu & Circuito protegido con interruptor termomagnetico de 10 A \\
Tomacorrientes & 2.5 mm2 Cu & Circuitos protegidos con interruptor termomagnetico de 16 A \\
Alimentación General & 10 mm2 Cu & Para la acometida desde el medidor de energía externa \\
Puesta a tierra & 6 mm2 Cu minimo & Conductor de protección conectado a barra de tierra \\
\bottomrule
\end{tabularx}
\end{table}
\end{landscape}

\section{Canalizaciones}

\begin{itemize}
  \item Tipo: tuberia PVC eléctrica, conduit u otra canalizacion permitida.
  \item Instalación: empotrada o superficial segun vivienda.
  \item Cajas: cajas octogonales para luminarias y cajas rectangulares para interruptores y tomacorrientes.
  \item Curvas y accesorios: compatibles con el diametro de tuberia.
\end{itemize}

\section{Tablero eléctrico}

El tablero debe permitir:
\begin{itemize}
  \item Interruptor general.
  \item Interruptor diferencial de protección de personas.
  \item Interruptores por circuito.
  \item Barra de neutro.
  \item Barra de tierra.
  \item Identificacion de circuitos.
  \item Espacio de reserva para ampliacion futura.
\end{itemize}

\section{Interruptores y protecciones}

\begin{itemize}
  \item Interruptores termomagneticos selecciónados segun corriente del circuito y conductor.
  \item Interruptor diferencial para protección de personas en tomacorrientes (30 mA).
  \item Coordinacion básica: la protección no debe superar la capacidad admisible del conductor ($I_b \le I_n \le I_z$).
\end{itemize}

\section{Tomacorrientes e interruptores}

\begin{itemize}
  \item Tomacorrientes con puesta a tierra en todos los ambientes.
  \item Alturas y ubicaciónes segun plano arquitectonico y criterio del curso.
  \item En baños y cocina se debe priorizar seguridad por humedad y uso de equipos.
\end{itemize}

\section{Luminarias}

\begin{itemize}
  \item Tipo: LED de adosar o empotrar.
  \item Potencia: registrar en cuadro de cargas.
  \item Ubicación: segun plano de alumbrado de cada piso.
\end{itemize}

\section{Sistema de puesta a tierra}

Componentes minimos:
\begin{itemize}
  \item Varilla o electrodo de puesta a tierra vertical (cobre electrolítico de 3/4" x 2.40 m).
  \item Conductor de cobre de protección de 6 mm2.
  \item Caja de registro de concreto H=0.30 m.
  \item Conector split-bolt de bronce.
  \item Barra de tierra en tableros.
\end{itemize}

\section{Cuadro de especificaciones}

\begin{landscape}
{\scriptsize
\setlength{\tabcolsep}{3pt}
\begin{longtable}{c L{3.0cm}L{1.0cm}r L{3.4cm}L{1.4cm}}
\caption{Especificaciones principales de materiales}\label{tab:especificaciones-materiales}\\
\toprule
\textbf{Item} & \textbf{Material} & \textbf{Und.} & \textbf{Cant.} & \textbf{Especificacion} & \textbf{Plano} \\
\midrule
\endfirsthead
\toprule
\textbf{Item} & \textbf{Material} & \textbf{Und.} & \textbf{Cant.} & \textbf{Especificacion} & \textbf{Plano} \\
\midrule
\endhead
1 & Conductor para alumbrado & m & 120 & Cobre 1.5 mm2 (fase, neutro) & IE-02/03/04 \\
2 & Conductor para tomacorrientes & m & 250 & Cobre 2.5 mm2 (fase, neutro y PE) & IE-02/03/04 \\
3 & Tuberia PVC-SEL & m & 180 & PVC eléctrica 3/4" de diámetro & IE-02/03/04 \\
4 & Tablero general/secundarios & und & 3 & Empotrados (TG-01, TD-01, TD-02) con barras N/T & IE-02/03/04 \\
5 & Interruptores termomagneticos & und & 7 & Interruptor General (1) y derivaciónes (6) & IE-05 \\
6 & Interruptor diferencial & und & 4 & General (1) y específicos de tomacorrientes (3) & IE-05 \\
7 & Puesta a tierra & glb & 1 & Varilla de cobre, registro y conductor & IE-06 \\
\bottomrule
\end{longtable}
}
\end{landscape}
